\section{K 近邻分类器}

\subsection{最近邻分类器的基本思想}

最近邻分类器（Nearest Neighbor）是最简单的数据驱动分类方法：

\begin{figure}[H]
\centering
\includegraphics[width=0.85\textwidth]{slides-images/slide-020.jpg}
\caption{最近邻分类器：\texttt{train} 仅记忆所有数据，\texttt{predict} 为每个测试样本找到最相似的训练样本\protect\footnotemark}
\end{figure}
\footnotetext{Slides 第21页。}

\begin{importantbox}{最近邻分类器的工作原理}
\begin{itemize}
    \item \textbf{训练阶段}：什么都不做，仅将所有训练数据和标签存入内存
    \item \textbf{预测阶段}：对于每个测试图像，计算它与\textbf{所有}训练图像的距离，返回距离最近的训练图像的标签
\end{itemize}
\end{importantbox}

\subsection{距离函数}

要比较两张图像的``相似度''，需要定义\textbf{距离函数}。最基本的两种距离度量是 L1 和 L2 距离。

\subsubsection{L1 距离（曼哈顿距离）}

\begin{figure}[H]
\centering
\includegraphics[width=0.85\textwidth]{slides-images/slide-022.jpg}
\caption{L1 距离计算示例：逐像素求差的绝对值，然后求和\protect\footnotemark}
\end{figure}
\footnotetext{Slides 第23页。}

$$d_1(I_1, I_2) = \sum_p |I_1^p - I_2^p|$$

其中：
\begin{itemize}
    \item $I_1, I_2$：两张图像
    \item $p$：遍历所有像素位置
    \item $I_1^p, I_2^p$：两张图像在位置 $p$ 的像素值
\end{itemize}

\subsubsection{L2 距离（欧氏距离）}

$$d_2(I_1, I_2) = \sqrt{\sum_p (I_1^p - I_2^p)^2}$$

\begin{figure}[H]
\centering
\includegraphics[width=0.85\textwidth]{slides-images/slide-030.jpg}
\caption{L1 距离 vs L2 距离的等距线可视化：L1 的等距线是菱形（正方形旋转 45$^\circ$），L2 的等距线是圆形\protect\footnotemark}
\end{figure}
\footnotetext{Slides 第31页。}

\begin{knowledgebox}{L1 与 L2 距离的核心区别}
\textbf{L1 距离}对坐标轴的选择\textbf{敏感}——如果旋转坐标系，L1 距离会发生变化。因此，当特征的每个维度都具有\textbf{明确的物理意义}时，L1 更合适。\\
\textbf{L2 距离}对坐标轴的选择\textbf{不敏感}——旋转坐标系不影响 L2 距离。因此，当特征维度之间没有特殊含义、较为\textbf{任意}时，L2 更合适。
\end{knowledgebox}

\subsection{最近邻分类器的实现}

\begin{figure}[H]
\centering
\includegraphics[width=0.95\textwidth]{slides-images/slide-025.jpg}
\caption{最近邻分类器的 Python 实现：仅需几行 NumPy 代码\protect\footnotemark}
\end{figure}
\footnotetext{Slides 第26页。}

\begin{lstlisting}[caption={最近邻分类器的核心实现},language=Python]
import numpy as np

class NearestNeighbor:
    def train(self, X, y):
        # X: N x D, y: N x 1
        self.Xtr = X
        self.ytr = y

    def predict(self, X):
        num_test = X.shape[0]
        Ypred = np.zeros(num_test, dtype=self.ytr.dtype)
        for i in range(num_test):
            # L1 distance
            distances = np.sum(np.abs(self.Xtr - X[i,:]), axis=1)
            min_index = np.argmin(distances)
            Ypred[i] = self.ytr[min_index]
        return Ypred
\end{lstlisting}

\subsection{计算复杂度分析}

\begin{table}[H]
\centering
\begin{tabular}{lcc}
\toprule
\textbf{操作} & \textbf{时间复杂度} & \textbf{说明} \\
\midrule
训练 & $O(1)$ & 仅存储数据，无计算 \\
预测（单样本） & $O(N)$ & 需要与所有 $N$ 个训练样本比较 \\
\bottomrule
\end{tabular}
\caption{最近邻分类器的计算复杂度}
\end{table}

\begin{warningbox}{预测慢、训练快——与实际需求恰好相反}
在实际应用中，我们希望训练可以慢（离线完成），但预测必须快（实时响应）。最近邻分类器恰好相反：训练 $O(1)$，预测 $O(N)$。这就像每次问 ChatGPT 一个问题，它都要遍历互联网上的所有答案——完全无法扩展。
\end{warningbox}

\subsection{从最近邻到 K 近邻}

\begin{figure}[H]
\centering
\includegraphics[width=0.95\textwidth]{slides-images/slide-028.jpg}
\caption{K 近邻的决策边界可视化：$K=1$（左）时存在噪声导致的错误区域，$K=3$（中）和 $K=5$（右）时决策边界更平滑\protect\footnotemark}
\end{figure}
\footnotetext{Slides 第29页。}

使用 $K=1$ 时，单个噪声点就会在其周围创建一个错误的决策区域。例如上图中，绿色区域中间有一个黄色点——这很可能是噪声，但 $K=1$ 会在其周围创建一大块黄色区域。

\begin{importantbox}{K 近邻的改进思路}
\textbf{K 近邻}（K-Nearest Neighbor, KNN）不再只看最近的一个邻居，而是找到最近的 $K$ 个邻居，通过\textbf{多数投票}（majority voting）决定分类结果。$K > 1$ 可以有效减少噪声和异常值的影响，使决策边界更加平滑。
\end{importantbox}

当 $K$ 较大时，可能出现票数相等的情况（白色区域），此时通常随机选择其中一个类别。这些``不确定区域''也为数据收集提供了有价值的指引——它们指示了需要更多样本的数据空间。

\begin{figure}[H]
\centering
\includegraphics[width=0.95\textwidth]{slides-images/slide-029.jpg}
\caption{不同 $K$ 值和距离函数的决策边界对比\protect\footnotemark}
\end{figure}
\footnotetext{Slides 第30页。}

\subsection{超参数与超参数调优}

KNN 中有两个需要选择的\textbf{超参数}（hyperparameters）：
\begin{itemize}
    \item $K$ 的值：使用多少个近邻
    \item 距离函数的选择：L1 还是 L2（或其他）
\end{itemize}

\begin{knowledgebox}{什么是超参数？}
超参数是\textbf{算法运行前}需要设定的参数，它们不是从数据中学习得到的，而是由人来选择的。超参数的最优值通常与具体的数据集和问题相关。
\end{knowledgebox}

\subsubsection{错误的超参数选择方法}

\begin{figure}[H]
\centering
\includegraphics[width=0.95\textwidth]{slides-images/slide-035.jpg}
\caption{超参数选择方法一：在训练数据上选择最优超参数——错误方法\protect\footnotemark}
\end{figure}
\footnotetext{Slides 第37页。}

\textbf{方法 1：在训练集上选择最优超参数}——这是最差的做法。对于 KNN，$K=1$ 在训练集上永远能达到 100\% 准确率（因为每个样本的最近邻就是自己），但这完全没有意义。

\textbf{方法 2：在测试集上选择最优超参数}——这也是错误的。虽然比方法 1 好一些，但它相当于用测试集``作弊''，无法保证模型在真正未见过的数据上的泛化能力。

\begin{warningbox}{绝对不要在测试集上调超参数}
测试集应该只在最终评估时使用\textbf{一次}。如果在测试集上反复调参，模型的性能评估将不可靠——你只知道它在这个特定测试集上表现如何，但不知道它能否泛化到其他数据。
\end{warningbox}

\subsubsection{正确的超参数选择方法}

\textbf{方法 3：训练集 + 验证集 + 测试集}

\begin{figure}[H]
\centering
\includegraphics[width=0.95\textwidth]{slides-images/slide-038.jpg}
\caption{正确做法：将数据分为训练集、验证集和测试集\protect\footnotemark}
\end{figure}
\footnotetext{Slides 第40页。}

从训练数据中划分出一部分作为\textbf{验证集}（validation set）。在训练集上训练模型，在验证集上评估不同超参数的效果，选择验证集上表现最好的超参数，最后在测试集上评估一次得到最终结果。

\textbf{方法 4：交叉验证（Cross-Validation）}

\begin{figure}[H]
\centering
\includegraphics[width=0.95\textwidth]{slides-images/slide-040.jpg}
\caption{$K$ 折交叉验证：将训练数据分成 $K$ 份（此处 $K=5$），每份轮流作为验证集，取平均准确率\protect\footnotemark}
\end{figure}
\footnotetext{Slides 第42页。}

\begin{importantbox}{交叉验证流程}
\begin{enumerate}
    \item 将训练数据分成 $K$ 等份（folds）
    \item 对于每一轮，用其中一份作为验证集，其余 $K-1$ 份作为训练集
    \item 对每种超参数设置，运行 $K$ 轮，取平均准确率
    \item 选择平均准确率最高的超参数
    \item 用选定的超参数在全部训练数据上重新训练，在测试集上评估
\end{enumerate}
交叉验证结果更可靠，但计算开销是单次验证的 $K$ 倍。在大规模深度学习中，由于计算成本过高，交叉验证较少使用，通常采用单一验证集的方法。
\end{importantbox}

\subsection{KNN 在 CIFAR-10 上的表现}

\begin{figure}[H]
\centering
\includegraphics[width=0.85\textwidth]{slides-images/slide-043.jpg}
\caption{CIFAR-10 数据集：10 个类别，50,000 张训练图像，10,000 张测试图像，每张 $32 \times 32 \times 3$\protect\footnotemark}
\end{figure}
\footnotetext{Slides 第45页。}

\begin{figure}[H]
\centering
\includegraphics[width=0.95\textwidth]{slides-images/slide-045.jpg}
\caption{KNN 在 CIFAR-10 上的最近邻检索结果：左侧为查询图像，右侧为按距离排序的 10 个最近邻\protect\footnotemark}
\end{figure}
\footnotetext{Slides 第47页。}

\begin{figure}[H]
\centering
\includegraphics[width=0.85\textwidth]{slides-images/slide-046.jpg}
\caption{五折交叉验证结果：$K=7$ 时获得最佳准确率约 28--29\%\protect\footnotemark}
\end{figure}
\footnotetext{Slides 第48页。}

在 CIFAR-10 上，KNN 最佳准确率约为 28\%。考虑到随机猜测的准确率为 10\%（10 个类别），KNN 确实在``做一些事情''，但还有很大提升空间。

观察检索结果可以发现，基于像素距离的匹配存在严重问题：颜色分布相似的图像会被认为``相近''，但它们在语义上可能完全不同。例如，一只绿色青蛙的最近邻可能是一只绿色背景的狗。

\begin{figure}[H]
\centering
\includegraphics[width=0.85\textwidth]{slides-images/slide-049.jpg}
\caption{像素距离的局限性：三张变换图像与原图的 L2 距离完全相同，但语义差异巨大\protect\footnotemark}
\end{figure}
\footnotetext{Slides 第51页。}

\begin{warningbox}{像素距离不适合图像分类}
将整张图像向右平移一个像素、改变颜色、或添加遮挡，都可能产生与原图\textbf{相同}的像素距离。像素级距离无法捕捉图像的语义信息，因此在实际中几乎不使用 KNN 进行图像分类。
\end{warningbox}

\subsection{KNN 小结}

\begin{figure}[H]
\centering
\includegraphics[width=0.85\textwidth]{slides-images/slide-050.jpg}
\caption{最近邻方法总结\protect\footnotemark}
\end{figure}
\footnotetext{Slides 第52页。}

\begin{itemize}
    \item KNN 是最简单的数据驱动分类方法，帮助我们理解了超参数和距离度量的概念
    \item 训练 $O(1)$，预测 $O(N)$——与实际需求恰好相反
    \item 像素级距离不适合图像分类
    \item 在实际中几乎不使用 KNN 处理图像，但它引出的概念（超参数、验证集、交叉验证）贯穿整个课程
\end{itemize}

\subsection{本章小结}

KNN 通过``记忆 + 查表''的方式工作，核心是距离函数的选择和 $K$ 值的设定。通过学习 KNN，我们掌握了数据驱动方法的基本范式、超参数调优的正确方法论（训练集/验证集/测试集的划分），以及交叉验证技术。同时，我们也认识到了基于原始像素距离的方法的根本局限性，这引出了更强大的参数化方法——线性分类器。

%% ========== Part 4: 线性分类器 ==========

